% !TEX program = xelatex

\documentclass[11pt,a4paper]{ctexart}

% ====== 字体 ======
\usepackage{fontspec}
\setmainfont{Latin Modern Roman}
\usepackage{unicode-math}
\setmathfont{Latin Modern Math}

% ====== 页面设置 ======
\usepackage[margin=2.5cm]{geometry}
\usepackage{setspace}
\onehalfspacing
\pagestyle{plain}

% ====== 数学宏包（三个经典功能，但用 unicode-math 替代了 amssymb） ======
\usepackage{mathtools}          % 必须放在 unicode-math 之前
\usepackage{amsmath}            % 已被 mathtools 自动载入，保留显式声明无妨
% \usepackage{amssymb}         % 删除！与 unicode-math 冲突，且 unicode-math 已提供所有符号

% ====== 三线表 + 数字对齐 ======
\usepackage{booktabs}
\usepackage{siunitx}
\usepackage{enumitem}          % 增强列表控制

% ====== 交叉引用与超链接 ======
\usepackage[colorlinks=true,
            linkcolor=blue,
            citecolor=green,
            urlcolor=cyan]{hyperref}
\usepackage[noabbrev,nameinlink]{cleveref}  % 必须最后加载

% ====== 自定义命令 ======
\newcommand{\dd}{\mathrm{d}}
\newcommand{\source}[1]{\par \hfill ——#1}

\title{志愿服务学习感想}
\author{王晨暄}
\date{\today}

\begin{document}
\maketitle

\section{一、万伏之下，匠心之上}

五月的主会场，窗外春光正好，室内却仿佛能听见高压电流的嗡鸣。王月鹏老师站在讲台上，
没有激昂的口号，只是平实地讲述着二十余年与“万伏”二字相伴的日子。那一刻，我忽然意识到，
真正的匠心从不在聚光灯下，而在那根悬于高空的银线上，在绝缘服被汗水浸透的褶皱里。

王老师分享的第一个案例，是一次深夜抢修。城市陷入黑暗，他穿上厚重的屏蔽服，
攀上数十米高的铁塔。万伏高压电场下，每一次呼吸都需精确计算。“不能快，不能慌，
每个动作都要像钟表齿轮般咬合。”他说得云淡风轻，我却在想象中感到指尖发麻。原来，
所谓“守光明”，守的不仅是城市的万家灯火，更是一个人对生命、对职责的敬畏边界。

真正触动我的，是他对一次技术革新的讲述。传统作业方式效率低、风险高，王老师带领团队历时三年，
将作业时间缩短了近一半。他说：“匠心不是重复，是在重复中看见改进的可能。”这让我反思，
我们这代人追逐“快”与“新”，却常常忘了，最深刻的创新往往源于最笨拙的坚守。
那些被汗水与时间淬炼出的“微创新”，才是大国工匠最动人的智慧。

讲座结束，那句“万伏高压守光明”仍在耳畔回响。我低头看自己的手——年轻、光滑，
尚未经历任何风霜。但王老师那双被绝缘手套磨出厚茧的手告诉我：青春践行使命，
不必惊天动地，只需在属于自己的“高压线”上，耐得住寂寞，顶得住压力，把平凡之事做到极致。

走出会场，阳光正好。我想，所谓大国工匠，不过是选择了最难走的路，却用最笨的方式，
走出了最远的征程。而我们这代人的使命，或许就是在各自的领域里，成为那根永不松懈的“高压线”——
沉默地承载，坚定地发光。


\section{二、守护青春底色，共筑健康中国}

近日，我有幸参加了学校举办的志愿大讲堂，聚焦HPV防治主题。
这场讲座不仅是一次知识的普及，更是一次深刻的青春启迪与国家
使命的共鸣。专家们的讲解，让我对“健康是青春最亮丽的底色”这
句话有了更真切、更厚重的理解。

健康，是青春最珍贵的基石。魏丽惠教授与王华庆教授以
详实的数据与生动的案例，深入浅出地阐述了HPV感染的普遍性与可防
可控性。这让我意识到，青春年华充满活力与可能，但一切梦想与奋
斗都离不开健康体魄的支撑。守护健康，绝非个人的私事，而是对自己
未来、对家庭幸福最基础的责任。讲座中提及的HPV与HIV（艾滋病病毒）
虽属不同病原体，但两者都警示我们，在充满活力的青年阶段，树立科学
的疾病预防观念、培养负责任的行为习惯至关重要。主动了解、积极预防
，是抵御包括HPV、HIV在内的各类健康风险的第一道防线，是对自己青春
最有力的护航。

国家的健康政策，为青春守护指明了方向与路径。张军院长关于
疫苗研发与应用的介绍，让我深刻感受到科技为健康守护带来的强大力量。
而这背后，正是国家“健康中国”战略的坚实支撑。党的二十大报告将健康中
国建设纳入中国式现代化布局，强调“健康第一”，这体现了国家对人民生命
健康、尤其是青年一代健康成长的高度重视。从将HPV疫苗纳入规划免疫的
积极考量，到广泛开展健康知识普及行动，一系列国家政策正在构建一个强
大的社会支持系统，为青年扫除健康认知的盲区，提供切实的保护工具。这
让我们看到，个人的健康守护并非孤军奋战，而是与国家发展同频共振。

作为新时代的青年，我们享受着国家提供的健康福祉，更应肩负起相
应的责任。我们不仅要成为自身健康的第一责任人，主动学习科学知识，接种
疫苗，定期筛查，养成良好生活习惯；更应成为健康中国的积极建设者和传播
者。正如本次讲座将健康守护与志愿服务深度融合所倡导的，我们可以把学到
的健康知识，特别是关于HPV、HIV等传染病的科学预防知识，通过志愿服务、
朋辈教育等形式传递给更多人，消除误解与歧视，营造关心健康、崇尚科学的社
会氛围。

这场讲座如同一盏明灯，照亮了健康之路。它让我坚信，将个人健康的微
观实践，融入国家健康政策的宏观布局，我们这一代青年必能以更加强健的
体魄、更加昂扬的姿态，在中国式现代化的壮阔征程中，绽放出最灿烂的青
春光芒，共同描绘健康中国的美好画卷。
\end{document}